\subsection{A fixed Lang cover counts prescribed primary components}
\label{sec:fixed-lang-density}

\paragraph{Finding A42 (an explicit finite-primary equidistribution refinement).}
Let $X/k$ be a smooth projective geometrically connected curve of
genus $g\ge1$ over $k=\mathbf F_q$, of characteristic $p$.
Let $B$ be a degree-one line bundle on $X$ defined over $k$, and put
\[
 J=\operatorname{Pic}^0(X),\qquad
 i_B(R)=\mathcal O_X(R)\otimes B^{-1}.
 \tag{PE.1}
\]
No rational effective representative of $B$, or rational point of
$X$, is assumed. Fix a finite nonempty set $S$ of rational primes,
which may contain $p$. For a torsion abelian group, a subscript $S$
denotes the sum of its primary components at primes in $S$. Put
\[
 G_S=J(k)_S,\qquad d=\#G_S,\qquad b=d(g-1)+1.
 \tag{PE.2}
\]
Let $\Phi$ be $q$-power Frobenius. For each $\ell\in S$, let
$h_\ell$ be the prime-to-$\ell$ part of the finite order of its
action on the \emph{geometric point group}
$J[\ell](\overline k)$; set $h_\ell=1$ if that group is zero.
Define
\[
 M=\operatorname{lcm}_{\ell\in S}(\ell h_\ell),\qquad
 \mathcal I=\{n\ge1:\gcd(n,M)=1\}.
 \tag{PE.3}
\]
For every $n\in\mathcal I$ there is equality of actual subgroups
of $J(\overline k)$,
\[
                    J(\mathbf F_{q^n})_S=G_S.
 \tag{PE.4}
\]
For each fixed $c\in G_S$ there is a smooth projective
geometrically connected curve $Y_c/k$, independent of $n$, and a
constant etale cover $Y_c\to X$ of degree $d$, such that
\[
 \begin{aligned}
 N_{c,n}&:=\#\{R\in X(\mathbf F_{q^n}):(i_B(R))_S=c\},\\
 dN_{c,n}&=\#Y_c(\mathbf F_{q^n}),\qquad n\in\mathcal I,\\
 \left|N_{c,n}-\frac{q^n+1}{d}\right|
             &\le\frac{2b}{d}q^{n/2}.
 \end{aligned}
 \tag{PE.5}
\]
In particular, along $n\in\mathcal I$ tending to infinity,
\[
                \frac{N_{c,n}}{\#X(\mathbf F_{q^n})}
                           \longrightarrow\frac1d.
 \tag{PE.6}
\]
Every prescribed component $c$ therefore occurs on infinitely
many points of the Abel curve, hence on a Zariski-dense subset.
The progression $n\equiv1\pmod M$ supplies arbitrarily large
admissible extension degrees without a theorem about primes.

\paragraph{Source credit and scope.}
The connected Lang quotient is the construction of Poonen,
\emph{Multiples of subvarieties in algebraic groups over finite
fields}, Lemma~5.1 (author PDF, p.~6). His proof of Theorem~1.4 on
the same page already controls primary subgroups by restricting
extension degrees; Lemmas~6.1--6.2, p.~7, give related finite-primary
extension bounds \cite{Poonen2005}. Milne's
\emph{Jacobian Varieties}, Proposition~9.1, p.~30, includes the
characteristic-primary etale part; its proof is completed on
pp.~31--32 \cite{MilneJV}. A12 already used the same connectedness
argument. We give the full fixed-cover, norm-scaling, and explicit
count refinement here, without a priority claim for the underlying
Lang or torsion-density argument.

\begin{proof}
\emph{1. Extension degrees that preserve every selected primary subgroup.}
For $n\ge1$ write
\[
 V_n(T)=1+T+\cdots+T^{n-1},\qquad
              \Phi^n-1=(\Phi-1)V_n(\Phi).
 \tag{PE.7}
\]
Fix $\ell\in S$ and $n\in\mathcal I$. On the finite
$\mathbf F_\ell$-vector space $J[\ell](\overline k)$,
$\Phi$ is invertible. Every eigenvalue $\lambda$ over
$\overline{\mathbf F}_\ell$ has order dividing $h_\ell$.
If $\lambda\ne1$, then $\lambda^n\ne1$ because
$\gcd(n,h_\ell)=1$, and
\[
 V_n(\lambda)=\frac{\lambda^n-1}{\lambda-1}\ne0.
\]
For $\lambda=1$ the value is $n\ne0$ in $\mathbf F_\ell$.
Thus $V_n(\Phi)$ is invertible on $J[\ell](\overline k)$,
also when $\Phi$ is not semisimple modulo $\ell$.

It is then injective on every finite group
$J[\ell^a](\overline k)$: a nonzero kernel would contain an
element of order $\ell$, contradicting the preceding conclusion.
It is therefore an automorphism on each such group. If
$x\in J(\overline k)$ is $\ell$-primary and $\Phi^n x=x$,
(PE.7) gives $V_n(\Phi)(\Phi-1)x=0$, hence $\Phi x=x$.
The converse is immediate. This proves (PE.4), prime by prime.
Only geometric torsion points are used; no assertion that the
group scheme $J[p]$ is etale is made.

This also proves the expected integral statement on the etale
$\ell$-adic Tate module: the same polynomial is an automorphism
because its reduction modulo $\ell$ is invertible. The finite-group
argument above suffices and avoids a choice of Tate-module basis.
If desired, one may replace $M$ by the explicit larger integer
\[
        M_0=\prod_{\ell\in S}
                   \left(\ell\prod_{j=1}^{2g}(\ell^j-1)\right).
 \tag{PE.8}
\]
Indeed the prime-to-$\ell$ part of the order of any element of
$\operatorname{GL}_r(\mathbf F_\ell)$, with $r\le2g$, divides
$\prod_{j=1}^{2g}(\ell^j-1)$. The characteristic-primary geometric
point group has rank at most $g$, so it satisfies this bound too.

\emph{2. One Lang quotient over the fixed ground field.}
Write $K=J(k)$ and let $H=K_{S^c}$ be its complementary primary
subgroup. The homomorphism $\Phi-1:J\to J$ is an etale isogeny:
its differential is $-1$ and its geometric kernel is the finite
constant group $K$. Quotienting by the constant subgroup $H$
gives
\[
 \begin{gathered}
 J\xrightarrow{\pi_H} A=J/H\xrightarrow{\psi}J,\\
       \psi\pi_H=\Phi-1,\qquad\ker\psi=K/H\simeq G_S.
 \end{gathered}
 \tag{PE.9}
\]
All maps and the indicated group identification are over $k$.
For $c\in G_S\subset J(k)$ form
\[
       Y_c=X\times_{i_B-c,\,J,\,\psi} A.
 \tag{PE.10}
\]
It is a constant etale $G_S$-torsor over $X$ of degree $d$.

\emph{3. Geometric connectedness, with no exclusion of $p$.}
Here is the universal-property proof, also showing that the
translation in (PE.10) causes no difficulty. Work over
$\overline k$ and choose $P_0\in X(\overline k)$.
If $Y_c$ were disconnected, its geometric monodromy would be a
proper subgroup $A_0\subset G_S$. Quotienting the isogeny by
$A_0$ would give a nontrivial intermediate etale isogeny
$\psi':A'\to J$ such that $i_B-c$ lifts to a morphism
$v:X\to A'$. The trivial monodromy of the quotient torsor is
exactly what supplies this lift. Subtracting $v(P_0)$ gives a
pointed lift of the usual Abel map $i_{P_0}$, since
\[
 \psi'\bigl(v(R)-v(P_0)\bigr)=i_B(R)-i_B(P_0)=[R-P_0].
\]
The Jacobian universal property supplies a homomorphism
$u:J\to A'$ with $\psi'u=\operatorname{id}_J$
\cite[Proposition~6.1, pp.~19--20]{MilneJV}.
The homomorphism $u$ is an isogeny, since it has finite kernel
and its image has the full dimension. Multiplicativity of degree
contradicts $\deg\psi'>1$. Thus $Y_c$ is geometrically connected.
It is smooth and projective, and unramified Riemann--Hurwitz gives
\[
                              g(Y_c)=d(g-1)+1=b.
 \tag{PE.11}
\]
This argument applies unchanged when $p$ divides $d$.

\emph{4. The exact norm and component detected by a split fibre.}
Let $k_n=\mathbf F_{q^n}$ and $R\in X(k_n)$, and put
$x=i_B(R)$. Choose $y\in J(\overline k)$ such that
$(\Phi-1)y=x-c$. The image $z=\pi_H(y)$ lies over $R$ in
(PE.10). Define
\[
 N_n=1+\Phi+\cdots+\Phi^{n-1}:J(k_n)\longrightarrow J(k).
\]
The target assertion follows from
$(\Phi-1)N_nx=(\Phi^n-1)x=0$. Direct telescoping gives
\[
 \begin{aligned}
 \Phi^n z-z
   &=\pi_H\bigl((\Phi^n-1)y\bigr)
     =\pi_H\bigl(N_nx-nc\bigr),\\
 \operatorname{pr}_{G_S}(N_nx-nc)
   &=n(x_S-c),\qquad n\in\mathcal I.
 \end{aligned}
 \tag{PE.12}
\]
For the second equality, (PE.4) places $x_S$ in $G_S$, on which
$\Phi$ acts trivially. The complementary component of $x$ has
order prime to every prime in $S$, so its image under $N_n$ has
the same property and projects to zero in $G_S$.
Multiplication by $n$ is an automorphism of $G_S$ because
$n$ is prime to every prime in $S$.
Consequently the fibre is split over $k_n$ if and only if
$x_S=c$. Since the deck group is constant, such a fibre has
exactly $d$ rational points, and every other fibre has none.
This proves the exact equality in (PE.5).

The first line of (PE.12) holds for every $n$; the second and the
split-fibre equivalence use the specified extension condition.
Thus no compatibility between different base-field Lang maps is
assumed. A single quotient (PE.9) and a fixed cover $Y_c/k$ are
being counted throughout. In particular, $n$ need not be one
modulo the exponent of $G_S$.

\emph{5. Counting, density, and descent.}
The Weil bound for the geometrically connected curve $Y_c$ gives
\[
       \bigl|\#Y_c(k_n)-(q^n+1)\bigr|\le2bq^{n/2}
\]
\cite[Theorem~11.1]{MilneJV}. Dividing by $d$ proves (PE.5).
Applying the same bound to $X$ proves (PE.6).
The numerator in the lower bound tends to infinity along
$n\equiv1\pmod M$, so the union of the counted point sets is
infinite. Since $i_B$ is an Abel embedding, it is Zariski dense
in its curve image. For a finite search, any admissible $n$ with
\[
                  q^n+1-2bq^{n/2}>d
 \tag{PE.13}
\]
gives at least two points in the desired primary coset. For
example $q^n>(2b+d)^2$ is sufficient. There is at most one point
with $i_B(R)=0$, so (PE.13) also guarantees a nonzero torsion
class when $c=0$.

Every object used in the construction is defined over the
specified finite field. If one first enlarges that field to
rationalize a desired finite torsion component or other input,
the theorem applies with that new $q$, $G_S$, and $M$. There is
no subsequent need to change the ground field or the Lang cover
as $n$ varies.
\end{proof}

\paragraph{Pure odd theta torsion and dense canonical-good points.}
Now let $C/\mathbf F_q$ be any smooth projective geometrically
connected genus-two curve in characteristic two. Set
$X=C^{(1)}$ and let $B$ be its degree-one locally exact
differential line bundle, so that $B^{\otimes2}\simeq\omega_X$.
The actual Raynaud theta support is
\[
                   \Theta=i_B(X)\subset J_X.
 \tag{PE.14}
\]
This line bundle and embedding are defined over $\mathbf F_q$;
no rational Weierstrass point is required. Taking $S=\{2\}$ and
$c=0$ in (PE.5) proves that theta points of odd order have
limiting proportion $1/\#J_X(\mathbf F_q)[2^\infty]$ along the
specified extension degrees. There are infinitely many nonzero
such points, for every $2$-rank and without a simplicity
assumption on the Jacobian. If the geometric $2$-torsion is
already rational, one may take every odd $n$; otherwise the
additional condition in (PE.3) excludes the finite set of
relevant unit-order prime divisors.

More generally, fix an odd prime $r$ with
$\operatorname{ord}_r(2)$ even; one may always choose $r=3$.
After a finite extension $k_0$ choose a nonzero
$c_r\in J_X[r](k_0)$. Apply the theorem over $k_0$ with
$S=\{2,r\}$ and $c=c_r$. It gives a Zariski-dense set of
points $R\in X(\overline k)$, with the explicit positive
proportion (PE.6) along the admissible extensions of $k_0$, for
which
\[
 \begin{aligned}
 (i_B(R))_2&=0,\qquad (i_B(R))_r=c_r\ne0,\\
 \kappa_R&=[K_X-2R]=-2i_B(R),\qquad
                  (\kappa_R)_r=-2c_r\ne0.
 \end{aligned}
 \tag{PE.15}
\]
At each such rational point, over the finite field in which it
was counted, A29 supplies an allowable pointed cyclic
$r$-power cover and separates its point section from those of
the other rational points. Frobenius twisting by the automorphism
of the finite ground field transports this statement, including
its rational points and pointed covers, to $C$.

\paragraph{What this does and does not settle.}
The characteristic-primary part of (PE.9) is an auxiliary
counting cover. It is \emph{not} a pro-prime-to-characteristic
neighbourhood and is never used as one. The allowable cover in
the last paragraph is the separate odd-degree construction of
A29. Pure odd theta torsion is now available quantitatively on
each fixed finite-field genus-two curve, and useful canonical
points are dense on each target. This does not force any image
of an arbitrarily prescribed source point, among the finitely
many images in a chosen dihedral cover, to be one of them;
A41 exhibits a fixed cover where all those images fail.
Nor does (PE.5) force the \emph{whole order} of a theta point
to be prime, a prime power, or an odd integer $m$ with
$\langle2,-1\rangle=(\mathbf Z/m)^\times$. Its unprescribed
primary components can contain additional primes. Consequently
the universal prescribed-cover choice and primitive-orbit
selection gates remain open; no previously open row is marked
closed by this density refinement.
